% 第7批：《大学物理简明教程》第7章静电场填空题7.2.1—7.2.12。
% 原书/PDF第186—187页；逐条独立研究性转述和推导，题图仍待二校。
\subsection{填空题：电场、通量、极化与电势}

\sourceproblem{7.2.1}{186}{186}
\problemtext 描述静电场的两个基本物理量及其一般定义式。
\workings
两个基本物理量为\textbf{电场强度$\bm E$}和\textbf{电势$\varphi$}。
电场强度定义为静止正试探电荷单位电量所受的力，
\[
\boxed{\bm E=\lim_{q_0\to0}\frac{\bm F}{q_0}}.
\]
以无穷远处电势为零的静电场中，
\[
\boxed{\varphi(P)=\frac{W_{P\to\infty}}{q_0}
=-\int_\infty^P\bm E\cdot d\bm l}.
\]
\crosscheck 电势零点在一般问题中可另选；不可把$\varphi=0$作为“无电场”的充分条件。

\sourceproblem{7.2.2}{186}{186}
\problemtext 静电场中若某一区域的电势在空间处处为同一常量，求该区域的电场强度。
\workings
静电场与电势关系为
\[
\boxed{\bm E=-\nabla\varphi}.
\]
若在一个开区域内$\varphi(\bm r)=\varphi_0$，各方向空间偏导数均为零，故
$\boxed{\bm E=\bm0}$。
\crosscheck 此结论是区域内处处电势不变；
单独某一点的电势等于零不能推出该点场强为零。

\sourceproblem{7.2.3}{186}{186}
\problemtext 点电荷$q$置于立方体中心，求穿过立方体任一面的电通量；
随后把该点电荷移到立方体外，求穿过整个立方体封闭表面的总电通量。
\workings
立方体总通量$\Phi=q/\varepsilon_0$。
当电荷位于中心时，由六个面在立方体对称变换下等价，
每面的\emph{有向}通量相同，
\[
\boxed{\Phi_{\mathrm{face}}=\frac{q}{6\varepsilon_0}}.
\]
电荷移至立方体外，封闭面内净电荷为零，故
\[
\boxed{\Phi_{\mathrm{closed}}=0}.
\]
\crosscheck 电荷位于外部后，每个面单独的通量不必为零；
只保证六个面的有向通量总和为零。

\sourceproblem{7.2.4}{186}{186}
\problemtext 绝缘薄橡胶球壳的总电荷$q$均匀分布于球面且保持不变。
球壳膨胀过程中球面越过距球心$r$处的固定观察点，问该点的电场强度大小如何变化。
\workings
在球面越过观察点之前，该点位于带电球壳外，
\[
E_{\mathrm{before}}=\frac{|q|}{4\pi\varepsilon_0r^2}.
\]
球壳半径变得大于$r$后，该点处于均匀带电球壳内部，
由球对称性与高斯定理
$\boxed{E_{\mathrm{after}}=0}$。
故场强由$\boxed{|q|/(4\pi\varepsilon_0r^2)}$变为$\boxed{0}$。
\crosscheck 观察点与球心距离$r$不变；
电荷只随球面均匀稀释，不能误把球面电荷密度降低直接代入外部点场强公式。

\sourceproblem{7.2.5}{186}{186}
\problemtext 一质子沿图7.2.5中$a\to c\to b$路径移动，电场力做功$8.0\times10^{-15}\,\mathrm J$。
求从$b$沿另一条路径返回$a$时电场力的功；
若$\varphi_a=0$，求$\varphi_b$。
\workings
静电场是保守场，电场力的功只取决于两个端点，
因此反向返回的功
\[
\boxed{W_{b\to a}=-8.0\times10^{-15}\,\mathrm J}.
\]
质子带电量$e=1.602176634\times10^{-19}$ C，
由$W_{a\to b}=e(\varphi_a-\varphi_b)$，
\[
\boxed{\varphi_b=\varphi_a-\frac{W_{a\to b}}e
\approx-4.99\times10^4\,\mathrm V}.
\]
\crosscheck 不能用折线或曲线的实际长度计算静电保守力功；
质子正电且场力做正功，终点电势更低。

\sourceproblem{7.2.6}{186--187}{186--187}
\problemtext 同心金属导体球壳，内壳净电荷$+q$、外壳净电荷$-2q$，
整个系统处于静电平衡。求外壳内表面、外表面的电荷量。
\workings
在外金属壳的材料内部取高斯面，
金属静电平衡有$\bm E=0$，故高斯面内净电荷为零：
$q+Q_{\mathrm{outer,in}}=0$，于是
\[
\boxed{Q_{\mathrm{outer,in}}=-q}.
\]
外球壳自身净电荷$-2q$，
\[
\boxed{Q_{\mathrm{outer,out}}=-2q-(-q)=-q}.
\]
\crosscheck 同心分布使面电荷均匀，但求总电荷不需要预先假定面密度均匀。

\sourceproblem{7.2.7}{187}{187}
\problemtext 将一负点电荷自无穷远移近一个起初不带电、\emph{未接地}的孤立金属导体。
求导体内部静电场强与导体电势相对初态的变化。
\workings
在任意静电平衡状态，导体内总电场$\bm E_{\mathrm{in}}=0$，
所以场强\textbf{不变}，仍为零。
但外部负电荷使导体表面自由电荷重新分布，
\emph{整块导体的电势}通常相对无穷远零点降低。
因此
\[
\boxed{\bm E_{\mathrm{in}}=\bm0}.
\]
导体电势相对初态\textbf{减小}。
\crosscheck 电势是全空间各源的叠加，不因导体内场强为零而必等于零；
孤立与接地是不同边界条件。

\sourceproblem{7.2.8}{187}{187}
\problemtext 一导体外侧填满均匀线性各向同性电介质（相对介电常数$\varepsilon_r$），
测得导体表面附近电场强度大小$E$，求表面自由电荷面密度的大小。
\workings
静电平衡时导体内电场零，界面法向位移矢量边界条件为
$\bm D_{\mathrm{out}}\cdot\bm n=\sigma_{\mathrm f}$。
在线性介质内$\bm D=\varepsilon_0\varepsilon_r\bm E$，
所以
\[
\boxed{|\sigma_{\mathrm f}|=\varepsilon_0\varepsilon_r E}.
\]
\crosscheck 若需要带符号的面密度，取$\sigma_f=\bm D\cdot\bm n$；
已知场强\emph{大小}时仅能求电荷面密度大小。

\sourceproblem{7.2.9}{187}{187}
\problemtext 平行板电容器板间完全充满均匀各向同性电介质$\varepsilon_r$，
忽略边缘效应。已知极板自由电荷面密度大小为$\sigma$，
求介质中的$\bm D$与$E$的大小。
\workings
无限大极板近似下边界条件给
\[
\boxed{D=\sigma},\qquad
\boxed{E=\frac{\sigma}{\varepsilon_0\varepsilon_r}}.
\]
方向由正极板指向负极板。
\crosscheck 这里$\sigma$表示\emph{自由}面电荷密度；
介质表面的极化束缚电荷已隐含在$\varepsilon_r$中，
不能再重复添加。

\sourceproblem{7.2.10}{187}{187}
\problemtext 两张无限大均匀带电平行平面位于真空，
面电荷密度分别为$+\sigma$和$+2\sigma$，板间距$d$。
求两平面之间的电场强度大小。
\workings
每一无限平面在任一侧产生大小$\sigma_i/(2\varepsilon_0)$的场，
方向远离正带电平面。
在两板之间两电场反向叠加，
\[
\boxed{E=\left|\frac{2\sigma}{2\varepsilon_0}
-\frac{\sigma}{2\varepsilon_0}\right|
=\frac{\sigma}{2\varepsilon_0}}.
\]
\crosscheck 该值与$d$无关；
不要将两板内场强错误相加为$3\sigma/(2\varepsilon_0)$。

\sourceproblem{7.2.11}{187}{187}
\problemtext 无外电场时分子正负电荷中心重合的电介质称什么；
施加外电场后分子正负电荷中心相对移动形成何种极化。
\workings
正负电荷中心原本重合的分子无永久电偶极矩，
属于\textbf{非极性电介质}。
外电场使电子云与原子核的中心相对错开，
产生感生电偶极矩，称为\textbf{位移极化}。
填空依次为：\textbf{非极性电介质；位移极化}。
\crosscheck 具有永久偶极矩的极性分子在外电场中还可产生转向极化；
“位移极化”不应误写成“转向极化”。

\sourceproblem{7.2.12}{187}{187}
\problemtext 半径$r_1=5\,\mathrm{cm}$、$r_2=20\,\mathrm{cm}$的同心均匀带电球面，
电荷分别$q_1=+3\times10^{-8}$ C、$q_2=-6\times10^{-8}$ C。
取无穷远零电势，求两球面之间电势为零的球面半径。
\workings
当$r_1<r<r_2$，内球面电势为$k_eq_1/r$，
外球面在其内部贡献常数$k_eq_2/r_2$，所以
\[
\varphi(r)=\frac{1}{4\pi\varepsilon_0}
\left(\frac{q_1}{r}+\frac{q_2}{r_2}\right).
\]
令$\varphi=0$得
\[
\boxed{r=-\frac{q_1}{q_2}r_2
=\frac{3}{6}\cdot20\,\mathrm{cm}
=10\,\mathrm{cm}}.
\]
\crosscheck $5<10<20$ cm落在题设区域中；
如果把外壳在壳内的电势也写成$k_eq_2/r$会得到错误答案。
